\section{Concrete source and table transformations}
Let $i$ be the source instruction incrementing register 1. Change the
constructor loop tail from
\[
 ((i,i),(i,i));\operatorname{continueTo}(2)
 \quad\hbox{to}\quad
 ((i,i),((i,i),\operatorname{continueTo}(3))).
\]
The continuation index changes with its syntactic depth. Both trees
execute four increments and jump to the same enclosing decrement.
In the named source DAG this changes \texttt{node62}. Remove the old
four-increment alias to native control 190: that control now participates
in continuation routing. The resolved instruction opcodes, register
arguments, and successor ordinals remain byte-for-byte equal to the
parent array. There are 3,671 instructions for GH and 3,668 for GHS;
305 source addresses change in each case.

Exactly the following five native transition rows change relative to
each parent. A row gives write symbol, movement, and next control.
\begin{center}
\begin{tabular}{rrll}
\toprule
Control & Read & Parent & New \\
\midrule
26 & 1 & $1\ R\ 42$ & $1\ R\ 91$ \\
101 & 0 & $0\ L\ 26$ & $0\ R\ 17$ \\
139 & 0 & $0\ R\ 190$ & $0\ R\ 99$ \\
139 & 1 & $0\ L\ 16$ & $1\ R\ 190$ \\
190 & 1 & $1\ R\ 99$ & $1\ L\ 73$ \\
\bottomrule
\end{tabular}
\end{center}
The first two rows shorten the two existing guarded decoder paths; the
last three implement the displayed source reassociation. The literal
checker verifies their combined effect on every source branch, so no
assumption that independently correct optimizations commute is needed.
The contextual folds are replayed simultaneously on the final table.
The count is $134+67-7=194$ for GH and $135+67-8=194$ for GHS.
Native start, source root, and halt controls are respectively 74, 84,
and 254. The delivered dense tables have start 0 and halt 194.
All 194 working states are graph-reachable from the dense start.

The exact dense hashes are
\begin{center}\small\ttfamily
\begin{tabular}{l}
PA194LCAX11GHF:\\
01c1b37ac3bbc51b2dfef22b8e42a6b3713ed720350269f819d82793e3f01bce\\[3pt]
PA194LCAX11GHSF:\\
6d03507fcc85825f81f7db5d02b0cb60e48f87c53be3dd0e4a8888d4922517a9.
\end{tabular}
\end{center}

\section{Literal certificates and their arithmetic meaning}
The construction-side symbolic replay checks 4,778 GH branch outcomes
and 4,774 GHS outcomes. The independent checker imports neither the
constructor nor arithmetic-box oracles. It derives four generic list
traversal lemmas from the actual rows, resolves the source control flow
by two separate algorithms, and checks every branch for nine arbitrary
natural registers, a nonnegative affine gap, any first spare, and any
nonempty finite spare list. It compares the entire finite-support tape
after blank startup, including blank exteriors.

The independent replay for the new GH table performs 1,105,240 symbolic
operations, and the GHS replay performs 1,109,186. These counts are
checker operations, not physical TM transitions. The maximum per-case
symbolic count is 385; every certified physical duration has constant
term at least 3. Exact physical startup lengths are 57,870 and 57,712.
The parent contracts retained from the supplied archive have startups
57,922 and 57,768. Their provenance is explicitly distinguished in the
coefficient receipt from the newly replayed contracts.

For arithmetic formalization represent a finite-support tape by a
zipper $(q,l,r)$, with the low bit of $r$ under the head and the low bit
of $l$ immediately to its left. After writing $w$, right and left moves
respectively send $(l,r)$ to
\[
 (2l+w,\lfloor r/2\rfloor),\qquad
 (\lfloor l/2\rfloor,4\lfloor r/2\rfloor+2w+(l\bmod2)).
\]
Table lookup, finite word operations, and bounded iteration are primitive
recursive. Absorb the halt control into a total transition function.
The symbolic executor uses the literal transition equation, quantifier-free
induction for homogeneous runs, and coded-list induction for the four
list traversals. The normalization identity
$C(A)1=1B(A)$, with $B(v)=0\,1^{v+1}$ and $C(v)=1\,0\,1^v$,
follows by induction on coded-list length. Exterior-word codes remain
free parameters. Canonical zipper equality replaces a quantified
assertion over all tape coordinates.

Every resulting local equation is a primitive-recursive frame equality.
All intermediate affine run lengths are coefficientwise nonnegative;
zero and successor decrement branches are checked separately. Finite
composition of these equations proves the local contracts in \ipra.
The executable checker is evidence for derivations using these sound
rules; Python is not asserted to be a proof-assistant kernel.

\section{An explicit relative theorem}
For any of the four tables $M$, let $H_M(t)$ mean that the totalized blank
run is halted after $t$ transitions. It is a decidable primitive-recursive
predicate, expressing halting by time $t$.
\begin{theorem}[Relative transfer]
For either corresponding parent $P$ and child $N$, the local contracts
and startups above yield
\[
 \ipra\vdash\forall t\,[H_P(t)\to H_N(\Theta(t))],\qquad
 \ipra\vdash\forall t\,[H_N(t)\to H_P(\Theta(t))],
\]
where
\[
 \boxed{\Theta(t)=57\,922+13(t+1)^2+689(t+1)
                   =13t^2+715t+58\,624.}
\]
Thus \ipra\ proves equivalence of the two halting sentences and, by
intuitionistic logic, equivalence of the two nonhalting sentences.
\end{theorem}
\begin{proof}
Augment a logical source configuration by its gap $g$, first spare $a_0$,
and finite spare tail $A$. Its physical encoding $E(s)$ is a primitive-
recursive function of these explicit arguments. The source counter
selects one of the finitely many valid instruction addresses. Every
source frame is physically nonhalting, even a frame whose next source
instruction is halt.

Let $h(s)$ mark the ready-to-execute source halt instruction. Let $S$ be
the augmented source step, with an absorbing logical halt extension;
let $s_n=S^n(s_0)$. For each nonhalting source state let $k(s)\ge1$ be the
exact physical duration and let $k_H(s)\ge1$ be the halt-site duration.
Extend these functions primitively recursively on unused inputs. With
startup length $b$, define
\[
 T(0)=b,\quad T(n+1)=T(n)+k(s_n),\quad
 B(n)=\prod_{j<n}(1-h(s_j)).
\]
The quantifier-free synchronization invariant is
\[
 B(n)=1\ \Longrightarrow\
 \bigl(\operatorname{Inv}(s_n)\land
 \operatorname{run}_M(T(n),\mathrm{blank})=E(s_n)
 \land T(n)\ge b+n\bigr).
\]
The complete startup proves its base. The equation
$B(n+1)=B(n)(1-h(s_n))$ supplies the nonhalt premise of the next local
contract. Exact iteration composition and duration positivity prove the
step. The frame target preserves its explicit auxiliary invariant.
Normal segments cannot hide a halt: in absorbing semantics an exact
endpoint equation at a nonhalting frame excludes an earlier halt.

Given a source halt witness $h(s_n)=1$, bounded first-hit minimization
returns its first index $m\le n$. Then $B(m)=1$, and the halt contract
gives a physical halt at $T(m)+k_H(s_m)$. Conversely, if $H_M(t)$ and
$B(t+1)=1$, synchronization produces a nonhalting root frame at time
$T(t+1)\ge b+t+1>t$, contradicting absorbing halting. Therefore the finite
Boolean prefix has a first hit $m\le t$. This is bounded minimization
with quantifier-free correctness induction; no Markov principle or
unbounded existential double-negation elimination is used.

The two implementations have identical resolved logical source graphs
and initial logical registers. Their auxiliary fields can differ, but
quantifier-free induction gives the same logical projection at every
source index. A source halt witness in either therefore transfers to
the other.

For a concrete uniform physical bound, the startup logical registers
are $(2,0,\ldots,0)$, the gap is zero, the first spare is 1, and the tail
consists of 25 ones. Its block width is 75. A finite coefficientwise
comparison of all four sets of contracts, tied by hashes to their tables,
gives
\[
 3\le k(s),k_H(s)\le
 542+8g+4\sum_{i=0}^{8}r_i+2a_0+2\operatorname{width}(A).
\]
A successor-decrement vector uses the decremented register value, which
can only reduce this upper bound. One source primitive increments at
most one register by one. Its gap rises by at most two; an increment
appends one spare value 1, and the other primitives retain the spare
list. Therefore, after $j$ normal source steps,
\[
 \sum_i r_i\le2+j,\quad g\le2j,\quad a_0=1,
 \quad\operatorname{width}(A)\le75+3j.
\]
All these are quantifier-free induction invariants. Substitution yields
$k(s_j),k_H(s_j)\le702+26j$. Hence $n$ normal source steps take at most
$b+13n^2+689n$. A first source halt at index $n$ takes at most
$b+13(n+1)^2+689(n+1)$. Since all four startups satisfy $b\le57\,922$
and extraction gave $n\le t$, the displayed $\Theta$ proves both
implications. Existential introduction gives halting equivalence;
contraposition of each implication, valid intuitionistically, gives
nonhalting equivalence.
\end{proof}

